% =============================================================================
% UnB Acessível — Documentação do Projeto
% Técnicas de Programação 2 — UnB 2026.1
%
% Para compilar:
%   xelatex documentacao.tex
%   xelatex documentacao.tex   (segunda vez para o sumário)
%
% Requer: TeX Live com xelatex e os pacotes listados abaixo.
% =============================================================================

\documentclass[12pt, a4paper]{article}
\special{papersize=210mm,297mm}

% --- Encoding e fontes (XeLaTeX) ---
\usepackage{fontspec}
\setmainfont{Liberation Sans}
\setmonofont{Liberation Mono}

% --- Idioma ---
\usepackage{polyglossia}
\setmainlanguage{brazil}

% --- Layout ---
\usepackage[top=2.5cm, bottom=2.5cm, left=2.5cm, right=2.5cm]{geometry}
\usepackage{setspace}
\onehalfspacing

% --- Cores ---
\usepackage{xcolor}
\definecolor{unbgreen}{HTML}{006633}
\definecolor{unbdark}{HTML}{1a1a2e}
\definecolor{codebg}{HTML}{f5f5f5}
\definecolor{linkblue}{HTML}{0066cc}
\definecolor{notebg}{HTML}{fff8e1}
\definecolor{noteborder}{HTML}{ffb300}

% --- Links ---
\usepackage{hyperref}
\hypersetup{
  colorlinks=true,
  linkcolor=unbdark,
  urlcolor=linkblue,
  pdftitle={UnB Acessível — Documentação do Projeto},
  pdfauthor={Equipe TP2 2026.1}
}

% --- Tabelas ---
\usepackage{array}
\usepackage{booktabs}
\usepackage{tabularx}
\usepackage{colortbl}

% --- Código ---
\usepackage{listings}
\lstset{
  backgroundcolor=\color{codebg},
  basicstyle=\ttfamily\small,
  breaklines=true,
  frame=single,
  rulecolor=\color{gray!30},
  xleftmargin=1em,
  framexleftmargin=0.5em,
  aboveskip=1em,
  belowskip=1em,
  literate={├}{{\textSFviii}}1 {└}{{\textSFii}}1 {│}{{\textSFx}}1 {─}{{\textSFi}}1
}

% --- Listas ---
\usepackage{enumitem}
\setlist[itemize]{leftmargin=1.5em, itemsep=0.3em}

% --- Cabeçalho e rodapé ---
\usepackage{fancyhdr}
\pagestyle{fancy}
\fancyhf{}
\fancyhead[L]{\small\color{gray}UnB Acessível}
\fancyhead[R]{\small\color{gray}TP2 — 2026.1}
\fancyfoot[C]{\small\color{gray}\thepage}
\renewcommand{\headrulewidth}{0.4pt}

% --- Caixas de nota ---
\usepackage{tcolorbox}
\tcbuselibrary{skins}
\newtcolorbox{notabox}{
  colback=notebg,
  colframe=noteborder,
  boxrule=0.5pt,
  left=8pt,
  right=8pt,
  top=6pt,
  bottom=6pt,
  arc=2pt,
  fonttitle=\bfseries,
  title=Nota
}

% --- Separador de seção ---
\newcommand{\partseparator}[1]{
  \clearpage
  \thispagestyle{empty}
  \vspace*{\fill}
  \begin{center}
    {\color{unbgreen}\rule{0.4\textwidth}{1pt}}\\[1.5em]
    {\Huge\bfseries\color{unbdark}#1}\\[1.5em]
    {\color{unbgreen}\rule{0.4\textwidth}{1pt}}
  \end{center}
  \vspace*{\fill}
  \clearpage
}

% =============================================================================
\begin{document}

% --- CAPA ---
\begin{titlepage}
  \thispagestyle{empty}
  \begin{center}
    \vspace*{2cm}

    {\color{unbgreen}\rule{\textwidth}{2pt}}

    \vspace{1.5cm}

    {\Huge\bfseries\color{unbdark}UnB Acessível}

    \vspace{0.8cm}

    {\LARGE\color{gray}Documentação do Projeto}

    \vspace{2cm}

    {\Large Técnicas de Programação 2}\\[0.5em]
    {\large Universidade de Brasília — 2026.1}

    \vspace{3cm}

    {\large Abril 2026}

    \vfill

    {\color{unbgreen}\rule{\textwidth}{2pt}}

    \vspace{1cm}

    {\small Hub de acessibilidade com FAQ inteligente e chatbot de IA\\para a comunidade universitária da UnB}
  \end{center}
\end{titlepage}

% --- SUMÁRIO ---
\tableofcontents
\clearpage

% =============================================================================
% PARTE I — README
% =============================================================================

\partseparator{Parte I — README do Projeto}

\section{UnB Acessível}

Hub de acessibilidade da Universidade de Brasília — um portal centralizado com FAQ inteligente e chatbot de IA para ajudar estudantes a encontrar informações sobre direitos, benefícios, apoio contra assédio e orientações institucionais.

\subsection{Sobre o Projeto}

A UnB disponibiliza diversas informações relevantes para seus estudantes — guias sobre assédio, benefícios socioeconômicos, orientações de fórum inicial e muito mais — mas essas informações estão espalhadas por diferentes sites, PDFs e setores. O \textbf{UnB Acessível} centraliza tudo isso em uma interface moderna e acessível, com um chatbot de IA treinado especificamente no conteúdo institucional da universidade.

Este projeto é desenvolvido como trabalho da disciplina de \textbf{Técnicas de Programação 2} do semestre 2026.1.

\subsection{Funcionalidades}

\begin{itemize}
  \item \textbf{FAQ Categorizado} — Perguntas frequentes organizadas por tema (assédio, benefícios, matrícula, saúde mental, etc.)
  \item \textbf{Chatbot com IA Local} — Modelo fine-tuned rodando no servidor para responder dúvidas em linguagem natural
  \item \textbf{Busca Inteligente} — Pesquisa por palavras-chave com ranking de relevância
  \item \textbf{Guias Passo a Passo} — Tutoriais visuais para processos burocráticos (como solicitar auxílio, como denunciar assédio)
  \item \textbf{Responsivo e Acessível} — Interface adaptada para mobile e compatível com leitores de tela (WCAG 2.1 AA)
  \item \textbf{Painel Administrativo} — Para a equipe manter o conteúdo atualizado sem precisar de deploy
\end{itemize}

\subsection{Stack Tecnológica (Sugestão Inicial)}

\begin{notabox}
A stack abaixo é apenas uma sugestão inicial e pode ser alterada conforme decisão da equipe durante o desenvolvimento.
\end{notabox}

\begin{center}
\begin{tabularx}{\textwidth}{>{\bfseries}l X}
  \toprule
  \rowcolor{unbgreen!10}
  Camada & Tecnologia \\
  \midrule
  Runtime & Bun \\
  Framework & React + TypeScript \\
  Build & Vite \\
  Estilização & Tailwind CSS + shadcn/ui \\
  IA / Chatbot & Modelo local fine-tuned (Ollama) \\
  Servidor & VPS \texttt{omarchy} \\
  \bottomrule
\end{tabularx}
\end{center}

\subsection{Estrutura do Projeto}

\begin{lstlisting}
unb-acessivel/
  src/
    components/     # Componentes React reutilizaveis
    pages/          # Paginas da aplicacao
    hooks/          # Custom hooks
    lib/            # Utilitarios e configuracoes
    services/       # Integracao com API e chatbot
    data/           # Conteudo do FAQ e guias
    types/          # Tipos TypeScript
  server/
    api/            # Endpoints da API
    ai/             # Integracao com modelo local
  public/           # Assets estaticos
  docs/             # Documentacao do projeto
    PRD.md          # Product Requirements Document
  README.md
\end{lstlisting}

\subsection{Requisitos}

\begin{itemize}
  \item Bun >= 1.0
  \item Node.js >= 20 (compatibilidade)
  \item Ollama (para o modelo de IA local)
\end{itemize}

\subsection{Instalação}

\begin{lstlisting}[language=bash]
# Clonar o repositorio
git clone https://github.com/fernando-augustop/trabalho-tp2-acessibilidade-unb-2026.1.git
cd trabalho-tp2-acessibilidade-unb-2026.1

# Instalar dependencias
bun install

# Rodar em desenvolvimento
bun dev
\end{lstlisting}

\subsection{Equipe}

Projeto desenvolvido por 11 alunos da disciplina de Técnicas de Programação 2 — UnB 2026.1.

\begin{center}
\begin{tabularx}{\textwidth}{X X}
  \toprule
  \rowcolor{unbgreen!10}
  \textbf{Aluno} & \textbf{GitHub} \\
  \midrule
  Fernando Augusto & \href{https://github.com/fernando-augustop}{@fernando-augustop} \\
  Gustavo Nascimento & \href{https://github.com/PavanelliGustavo}{@PavanelliGustavo} \\
  Eduardo Rocha & \href{https://github.com/eduardofgc}{@eduardofgc} \\
  Samara Gomes & \href{https://github.com/samaragomess}{@samaragomess} \\
  Auguto Faller & \href{https://github.com/tosgual}{@tosgual} \\
  \emph{A preencher} & \\
  \emph{A preencher} & \\
  \emph{A preencher} & \\
  \emph{A preencher} & \\
  \emph{A preencher} & \\
  \emph{A preencher} & \\
  \bottomrule
\end{tabularx}
\end{center}

\subsection{Licença}

Este projeto é de uso acadêmico, desenvolvido para a Universidade de Brasília.

% =============================================================================
% PARTE II — PRD
% =============================================================================

\partseparator{Parte II — Product Requirements Document}

\section{PRD — UnB Acessível}

\begin{tabularx}{\textwidth}{>{\bfseries}l X}
  Produto: & UnB Acessível — Hub de Acessibilidade da UnB \\
  Disciplina: & Técnicas de Programação 2 — 2026.1 \\
  Equipe: & 11 alunos \\
  Data: & 07 de abril de 2026 \\
  Versão: & 1.0 \\
\end{tabularx}

\subsection{Visão Geral}

\subsubsection{Problema}

A Universidade de Brasília oferece diversos recursos de apoio aos estudantes — desde guias sobre como agir em casos de assédio até procedimentos para solicitar benefícios socioeconômicos, orientações de fórum inicial e informações sobre saúde mental. Porém, esse conteúdo está fragmentado entre diferentes sites institucionais, PDFs dispersos, redes sociais e setores administrativos. O estudante que precisa de ajuda muitas vezes não sabe onde procurar ou desiste diante da burocracia.

\subsubsection{Solução}

O \textbf{UnB Acessível} é um hub web centralizado que reúne todas essas informações em um único lugar, com interface moderna, busca inteligente e um chatbot de IA treinado especificamente no conteúdo da universidade. Não é um portal de notícias — é um \textbf{FAQ moderno e interativo} que responde às dúvidas dos estudantes de forma direta e acessível.

\subsubsection{Objetivos}

\begin{enumerate}
  \item Centralizar informações institucionais de apoio ao estudante em um único portal
  \item Permitir que estudantes encontrem respostas rapidamente via busca ou chatbot
  \item Reduzir a barreira de acesso a informações sobre direitos e benefícios
  \item Garantir acessibilidade (WCAG 2.1 AA) para todos os usuários
  \item Manter o conteúdo atualizável sem necessidade de deploys
\end{enumerate}

\subsection{Público-Alvo}

\subsubsection{Usuários Primários}

\begin{itemize}
  \item \textbf{Estudantes de graduação e pós-graduação da UnB} — especialmente calouros e alunos em situação de vulnerabilidade socioeconômica
  \item \textbf{Estudantes que sofreram ou presenciaram assédio} — buscando orientação sobre como proceder
  \item \textbf{Estudantes com dúvidas burocráticas} — matrícula, trancamento, aproveitamento de créditos, etc.
\end{itemize}

\subsubsection{Usuários Secundários}

\begin{itemize}
  \item \textbf{Servidores e docentes} — para direcionar alunos ao recurso correto
  \item \textbf{Administradores de conteúdo} — equipe responsável por manter as informações atualizadas
\end{itemize}

\subsection{Funcionalidades}

\subsubsection{MVP (Escopo do Trabalho)}

\paragraph{F1 — FAQ Categorizado}

\begin{itemize}
  \item Página inicial com categorias temáticas (cards ou grid)
  \item Categorias planejadas:
    \begin{itemize}
      \item Assédio e Violência
      \item Benefícios e Auxílios Socioeconômicos
      \item Saúde Mental e Apoio Psicológico
      \item Matrícula e Vida Acadêmica
      \item Fórum Inicial e Recepção de Calouros
      \item Acessibilidade e Inclusão (PCD)
      \item Moradia e Restaurante Universitário
      \item Estágio e Oportunidades
    \end{itemize}
  \item Cada categoria contém perguntas e respostas expansíveis (accordion)
  \item Conteúdo armazenado em arquivos estruturados (JSON/MDX) para facilitar edição
\end{itemize}

\paragraph{F2 — Chatbot com IA Local}

\begin{itemize}
  \item Widget de chat flutuante acessível em todas as páginas
  \item Modelo de linguagem fine-tuned rodando localmente via Ollama
  \item Treinado com o conteúdo do FAQ e documentos institucionais da UnB
  \item Respostas em português brasileiro, linguagem simples e direta
  \item Indicação de fontes/links quando aplicável
  \item Fallback para ``não sei responder, entre em contato com [setor]''
\end{itemize}

\paragraph{F3 — Busca Inteligente}

\begin{itemize}
  \item Campo de busca no topo da página
  \item Busca full-text no conteúdo do FAQ
  \item Sugestões de perguntas enquanto o usuário digita (autocomplete)
  \item Ranking de relevância nos resultados
\end{itemize}

\paragraph{F4 — Guias Passo a Passo}

\begin{itemize}
  \item Páginas dedicadas para processos burocráticos comuns
  \item Formato de tutorial com etapas numeradas
  \item Capturas de tela dos sistemas da UnB quando aplicável
  \item Links diretos para os formulários e sistemas necessários
\end{itemize}

\paragraph{F5 — Design Responsivo e Acessível}

\begin{itemize}
  \item Mobile-first
  \item Conformidade com WCAG 2.1 nível AA
  \item Suporte a navegação por teclado
  \item Compatibilidade com leitores de tela
  \item Contraste adequado e tamanhos de fonte ajustáveis
  \item Modo escuro
\end{itemize}

\subsubsection{Futuro (Pós-Disciplina)}

\begin{itemize}
  \item Painel administrativo para CRUD de conteúdo
  \item Autenticação via SSO da UnB
  \item Sistema de feedback por pergunta (``essa resposta ajudou?'')
  \item Analytics de perguntas mais acessadas
  \item Notificações sobre prazos importantes (matrícula, auxílios)
  \item Versão PWA para acesso offline
\end{itemize}

\subsection{Arquitetura Técnica}

\begin{notabox}
A stack e arquitetura descritas abaixo são uma sugestão inicial. A equipe pode optar por outras tecnologias conforme preferência e experiência do grupo.
\end{notabox}

\subsubsection{Frontend}

\begin{lstlisting}
React + TypeScript + Vite
  Estilizacao: Tailwind CSS + shadcn/ui
  Roteamento: React Router
  Estado: React Context / Zustand (se necessario)
  Conteudo: Arquivos MDX ou JSON estaticos
\end{lstlisting}

\subsubsection{Backend / API}

\begin{lstlisting}
Bun (servidor)
  API REST para o chatbot
  Endpoint de busca
  Proxy para o modelo Ollama
\end{lstlisting}

\subsubsection{IA / Chatbot}

\begin{lstlisting}
Ollama (servidor local na VPS)
  Modelo base: Llama 3 ou Mistral (a definir)
  Fine-tuning com dados institucionais da UnB
  RAG (Retrieval-Augmented Generation) com embeddings do FAQ
  Contexto limitado ao dominio da UnB
\end{lstlisting}

\subsubsection{Infraestrutura}

\begin{lstlisting}
VPS omarchy
  Aplicacao servida via Bun/Caddy
  Ollama rodando o modelo de IA
  Deploy via Git (push to deploy) ou CI simples
  HTTPS via Caddy/Let's Encrypt
\end{lstlisting}

\subsubsection{Diagrama Simplificado}

\begin{lstlisting}
  Browser       -->   Bun Server   -->    Ollama
  (React SPA)   <--   (API REST)   <--   (LLM Local)
                          |
                      Conteudo
                      (JSON/MDX)
\end{lstlisting}

\subsection{Conteúdo Inicial}

\subsubsection{Fontes de Dados}

O conteúdo será compilado a partir de:

\begin{itemize}
  \item Site institucional da UnB (unb.br)
  \item Decanato de Assuntos Comunitários (DAC)
  \item Decanato de Ensino de Graduação (DEG)
  \item Ouvidoria da UnB
  \item Centro de Atendimento e Estudos Psicológicos (CAEP)
  \item Diretoria de Diversidade (DIV)
  \item Secretaria de Administração Acadêmica (SAA)
  \item Guias e cartilhas já publicados pela universidade
\end{itemize}

\subsubsection{Formato do Conteúdo}

Cada entrada do FAQ será um arquivo estruturado:

\begin{lstlisting}
{
  "id": "assedio-001",
  "categoria": "assedio-e-violencia",
  "pergunta": "Sofri assedio na universidade. O que devo fazer?",
  "resposta": "Voce pode procurar a Ouvidoria da UnB...",
  "links": [
    { "texto": "Ouvidoria da UnB", "url": "https://ouvidoria.unb.br" }
  ],
  "tags": ["assedio", "denuncia", "ouvidoria"],
  "atualizado_em": "2026-04-07"
}
\end{lstlisting}

\subsection{Requisitos Não-Funcionais}

\begin{center}
\begin{tabularx}{\textwidth}{>{\bfseries}l X}
  \toprule
  \rowcolor{unbgreen!10}
  Requisito & Meta \\
  \midrule
  Tempo de carregamento & < 2s (First Contentful Paint) \\
  Tempo de resposta do chatbot & < 5s \\
  Disponibilidade & 99\% (VPS dedicada) \\
  Acessibilidade & WCAG 2.1 AA \\
  Compatibilidade & Chrome, Firefox, Safari, Edge (últimas 2 versões) \\
  Mobile & Responsivo, touch-friendly \\
  SEO & Meta tags, sitemap, conteúdo indexável \\
  \bottomrule
\end{tabularx}
\end{center}

\subsection{Organização da Equipe}

\subsubsection{Divisão por Squads}

Com 11 alunos, a sugestão de divisão é:

\begin{center}
\begin{tabularx}{\textwidth}{>{\bfseries}l X l}
  \toprule
  \rowcolor{unbgreen!10}
  Squad & Responsabilidade & Tamanho \\
  \midrule
  Frontend Core & Layout, componentes, roteamento, responsividade & 3 alunos \\
  FAQ e Conteúdo & Estrutura do FAQ, busca, guias, curadoria de conteúdo & 3 alunos \\
  Chatbot e IA & Integração Ollama, RAG, fine-tuning, API do chat & 3 alunos \\
  Infra e DevOps & Deploy, VPS, CI/CD, domínio, HTTPS & 2 alunos \\
  \bottomrule
\end{tabularx}
\end{center}

\subsubsection{Workflow de Desenvolvimento}

\begin{itemize}
  \item Branch principal: \texttt{main} (protegida)
  \item Feature branches: \texttt{feat/<nome-da-feature>}
  \item Pull Requests obrigatórios com pelo menos 1 review
  \item Commits em português, formato: \texttt{tipo: descrição} (ex: \texttt{feat: adiciona página de FAQ})
\end{itemize}

\subsection{Cronograma Estimado}

\begin{center}
\begin{tabularx}{\textwidth}{>{\bfseries}l X}
  \toprule
  \rowcolor{unbgreen!10}
  Fase & Entregas \\
  \midrule
  Fase 1 — Setup & Repo, stack configurada, deploy básico, estrutura de pastas \\
  Fase 2 — Core & FAQ categorizado, layout principal, busca básica \\
  Fase 3 — IA & Chatbot funcional, RAG com conteúdo do FAQ \\
  Fase 4 — Conteúdo & Curadoria completa, guias passo a passo \\
  Fase 5 — Polish & Acessibilidade, testes, performance, documentação \\
  Entrega Final & Apresentação, demo ao vivo, documentação completa \\
  \bottomrule
\end{tabularx}
\end{center}

\subsection{Riscos e Mitigações}

\begin{center}
\begin{tabularx}{\textwidth}{l l X}
  \toprule
  \rowcolor{unbgreen!10}
  \textbf{Risco} & \textbf{Impacto} & \textbf{Mitigação} \\
  \midrule
  IA pesada para a VPS & Chatbot lento & Modelo menor (Phi-3, TinyLlama) ou RAG sem fine-tuning \\
  Conteúdo desatualizado & Respostas incorretas & Data de atualização e links para fonte oficial \\
  Escopo grande demais & Entrega incompleta & Priorizar MVP (FAQ + busca), chatbot como bônus \\
  Problemas de acessibilidade & Exclusão de usuários & Testes com Lighthouse e axe-core desde o início \\
  \bottomrule
\end{tabularx}
\end{center}

\subsection{Métricas de Sucesso}

\begin{itemize}
  \item FAQ com pelo menos 50 perguntas e respostas validadas
  \item Chatbot respondendo corretamente > 80\% das perguntas de teste
  \item Score de acessibilidade > 90 no Lighthouse
  \item Tempo de carregamento < 2s
  \item Deploy funcional e acessível via URL pública
\end{itemize}

\subsection{Referências}

\begin{itemize}
  \item \href{https://www.w3.org/TR/WCAG21/}{WCAG 2.1}
  \item \href{https://ui.shadcn.com}{shadcn/ui}
  \item \href{https://ollama.ai}{Ollama}
  \item \href{https://vitejs.dev}{Vite}
  \item \href{https://bun.sh}{Bun}
  \item \href{https://unb.br}{UnB — Universidade de Brasília}
\end{itemize}

\end{document}
